%!TEX root = ../thesis.tex
% ===================================================================
% Chapter 4: Implementation (replaces the stubs in chapters/implementation.tex)
% Uses the same tikzset as design_chapter.tex. If design.tex already defines
% it (it does, as of 11 Oct), delete the \tikzset block below to avoid a
% duplicate-definition warning; otherwise keep it, since this chapter also
% draws two figures.
% ===================================================================
\mychapter{Implementation}{Implementation}{}
\label{chap:implementation}

\tikzset{
  blk/.style={draw, rounded corners=2pt, align=center, font=\footnotesize,
              minimum width=2.6cm, minimum height=0.95cm, inner sep=3pt},
  mdl/.style={blk, fill=black!8},
  hw/.style={blk, dashed},
  st/.style={draw, rounded corners=2pt, align=center, font=\footnotesize,
             minimum width=2.7cm, minimum height=0.9cm},
  arr/.style={-{Stealth[length=2mm]}, thick},
  lbl/.style={font=\scriptsize, fill=white, inner sep=1pt}
}

Chapter~\ref{chap:design} said what each part of the pipeline does and why.
This chapter says how each part is actually built: the libraries, the process
layout, the message formats, and the places where a design decision turned
into a specific piece of code. I keep the numbers here that
Chapter~\ref{chap:design} left out on purpose, such as port numbers, file
names and exact thresholds, since those belong to how the system is built
rather than why it is shaped the way it is. Where Chapter~\ref{chap:design}
described a behaviour only as a concept, such as mid-mission revision or the
guidance loop that follows a blocked step, this chapter is also where that
concept gets grounded in the function that actually carries it out.

%% ===================================================================
\mysection{Speech input and output}{Speech Input and Output}
\label{sec:speech}

\mysubsection{Speech-to-text}{Speech-to-text}
\label{sec:stt}

Transcription runs through \texttt{faster-whisper}, a CTranslate2 reimplementation
of Whisper, using the \texttt{digiphyte/\allowbreak fluister-turbo} model. I picked that model
over stock Whisper because it is trained to include South African English and
Afrikaans accents, which matters for a project the supervisor will test in person.
The model runs on the CPU with \texttt{compute\_type="int8"} and 8 threads. I kept
it off the GPU on purpose, since the planner and the vision model already share
the 12\,GB of video memory, and transcription is short enough that CPU inference
does not become the bottleneck in a turn.

\texttt{transcribe\_audio} takes either raw bytes (an encoded clip from the
dashboard or the Pi client, decoded through FFmpeg) or a NumPy array of raw
float32 samples, since the capture path that reads straight from a microphone has
no container format to decode. The function is called once per utterance rather
than as a stream. Whisper pads every clip to 30\,s internally, so chunking a long
recording into pieces would not reduce latency, only repeat the padding cost.

One behaviour was worth adding after seeing it happen on real recordings.
\texttt{fluister-turbo} occasionally transcribes silence or background noise as
one of a handful of fixed captioning artefacts, most often ``thank you so much
for watching''. I did not expect this and found it by reading failed-command
logs where the ambiguity check was asked to resolve a command that was never
actually said. \texttt{transcribe\_audio} now checks the lowercased, punctuation-stripped
output against a short set of these phrases and returns an empty string
when it matches, so the rest of the pipeline treats it the same as a transcription
that genuinely failed, rather than as a real command.

\mysubsection{Speech output: Piper}{Speech Output}
\label{ssec:tts}

Spoken output goes through Piper, a local neural TTS engine, using the
\texttt{en\_US-lessac-medium} voice. \texttt{tts.py} mirrors \texttt{stt.py}'s
cached-model pattern: \texttt{get\_voice} loads the \texttt{PiperVoice} once
from \texttt{PIPER\_VOICE\_MODEL} and keeps it in a module-level variable, and
\texttt{synthesize\_speech} writes the result into an in-memory \texttt{wave}
buffer rather than a temporary file, since \texttt{synthesize\_wav} needs an
actual \texttt{Wave\_write} object to call \texttt{setframerate}\allowbreak\texttt{/writeframes}
on. Piper's voices are not fetched the way \texttt{faster-whisper}'s model is;
I had to run \texttt{python -m piper.\allowbreak download\_voices en\_US-\allowbreak lessac-medium} once
by hand and point the environment variable at the resulting \texttt{.onnx} file,
with its paired \texttt{.onnx.json} config sitting alongside it, which is how
Piper expects a voice to be packaged.

There is no Afrikaans voice in the \texttt{rhasspy/\allowbreak piper-voices} set, which I
checked rather than assumed. This turned out not to matter: every \texttt{\_speak}
call site speaks a hard-coded English template string regardless of the
session's detected language, so only the speech-to-text side needs to branch on
language at all. Speaking back in the user's own language would mean
localising those template strings, which is a separate piece of work from
picking a synthesis engine, and I have not done it.

Synthesis itself happens on a worker thread, through
\texttt{asyncio.to\_thread}, inside the same function,
\texttt{mission.\_speak}, that builds the \texttt{speak} event the Pi client
plays back. This matters because \texttt{\_speak} runs inside the step loop,
which shares an event loop with \texttt{ingest\_frame}; synthesis is sub-second
but still a blocking call, and perception should not stall mid-mission for the
sake of a voice line. If synthesis itself throws, the exception is caught and
logged, and the event is sent without an \texttt{audio} field: the mission
degrades to text-only rather than stopping, the same soft-fail stance the
pre-departure dialogue's own \texttt{\_speak} in \texttt{pipeline.py} already
takes.

%% ===================================================================
\mysection{Model inference layer}{Model Inference Layer}
\label{sec:inference}

\mysubsection{Planner}{Planner}
\label{sec:planner}

The planner is one function, \texttt{generate\_plan}, behind a provider interface
that Section~\ref{sec:vlm} shares with the vision model. \texttt{InferenceProvider}
is an abstract class with one method, \texttt{complete}, that takes a list of
chat messages and an optional image and returns a \texttt{Completion}: the text,
which provider and model answered, the latency, and token counts and cost when
the provider reports them. \texttt{OllamaProvider} and
\texttt{OpenAICompatibleProvider} both implement it, the second covering OpenAI
and DeepSeek, since both speak the same \texttt{/chat/completions} schema.
Swapping the planner or the vision model is then a matter of an environment
variable (\texttt{PLANNER\_PROVIDER}, \texttt{VLM\_PROVIDER}) and not a code
change, which is what the cloud comparison in
Section~\ref{sec:cloud_comparison} depends on.

\begin{figure}[htbp]
  \centering
  \begin{tikzpicture}[node distance=0.8cm]
    \node[blk] (pl) at (0,1.2)    {\texttt{planner.py}};
    \node[blk] (vl) at (0,-1.2)   {\texttt{vlm.py}};
    \node[mdl] (if) at (4.0,0)    {\texttt{InferenceProvider}\\\texttt{.complete()}};
    \node[mdl] (ol) at (8.5,1.2)  {\texttt{OllamaProvider}\\(local)};
    \node[mdl] (oa) at (8.5,-1.2) {\texttt{OpenAICompatibleProvider}\\(OpenAI, DeepSeek)};
    \draw[arr] (pl)--(if);
    \draw[arr] (vl)--(if);
    \draw[arr] (if)--(ol);
    \draw[arr] (if)--(oa);
  \end{tikzpicture}
  \caption{Both roles call the same interface. \texttt{get\_provider(role)} resolves \texttt{PLANNER\_PROVIDER} or \texttt{VLM\_PROVIDER} to a concrete class and caches the instance.}
  \label{fig:providers}
\end{figure}

\texttt{get\_provider(role)} resolves the provider name and the model tag from
configuration, with the model name itself looked up per provider
(\texttt{PLANNER\_MODEL\_OLLAMA}, \texttt{PLANNER\_MODEL\_OPENAI}, and so on), so
switching providers does not require re-typing a model name that was already
correct. A DeepSeek instance is constructed with \texttt{supports\_vision=False},
and \texttt{get\_provider} refuses to hand a text-only provider to the VLM role,
since DeepSeek's hosted models cannot see an image at all. Every construction is
cached by (role, provider, model), so a long-running process reuses one client
rather than reconnecting on every call.

\texttt{OllamaProvider} sends \texttt{num\_ctx} (8192 by default), an optional
\texttt{repeat\_penalty}, and \texttt{format="json"} for every planner and VLM
call, so the model is forced to emit JSON rather than prose. Two failures are
retried once before being allowed to fail a mission: a runner crash during a cold
CUDA load on \texttt{gemma4:e4b}, which happened during testing on 5 October, and
Ollama's own repeat-limit abort, which happens when a small model loops on the
same tokens. Both are transient, so a retry after a short sleep usually succeeds,
and it is cheaper than treating every such failure as a dead provider.

\texttt{build\_plan\_prompt} assembles the prompt from the active planner
profile (Section~\ref{sec:vocabulary}), the room digest when one is loaded, and
the conditional guidance block when the command contains a condition word
(Section~\ref{sec:conditional}). It is kept as a pure function, separate from
\texttt{generate\_plan} itself, so a small tool,
\texttt{tools/snapshot\_prompt.py}, can print the exact prompt a command would
produce without calling a model, which is how I checked the prompt actually
changes shape for a conditional command.

\texttt{generate\_plan} parses the model's JSON reply defensively rather than
trusting it. A smaller model sometimes writes a step as a bare string instead of
an object, or nests the list oddly; any step that is not a well-formed object is
dropped before anything downstream calls \texttt{.get} on it, and a note is
appended to the plan recording that steps were dropped. \texttt{\_promote\_moves\_to\_approach}
then rewrites a \texttt{move} step whose target names an object into an
\texttt{approach} step, the one repair mentioned in
Section~\ref{sec:vocabulary}.

Failure recovery (Section~\ref{ssec:guidance_loop}) is a second prompt builder,
\texttt{build\_replan\_prompt}, and a second entry point,
\texttt{replan\_from\_state}, rather than a parameter on \texttt{generate\_plan}.
I kept them separate because the two prompts carry different context: a first
plan is built from the command and the room digest, while a replan is built from
the command, the steps already done, the step that failed and why, and the
guidance conversation so far. \texttt{validate\_steps} runs after both, but with
a different action list depending on the call, since a replan in free-roam mode
has no \texttt{follow\_line} to offer (Section~\ref{sec:vocabulary}).

\mysubsection{Vision-language verification}{Vision-language Verification}
\label{sec:vlm}

Every call into the vision model goes through one function, \texttt{\_ask}, which
takes a named stage (for logging), a prompt, an optional image and a salvage
flag. It resolves the VLM provider, sends the prompt with \texttt{json\_mode=True},
and parses the reply. A reply that fails to parse as JSON is not immediately
treated as an error: when \texttt{salvage=True}, a short recovery routine,
\texttt{\_salvage\_truncated\_json}, is tried first. It walks the text tracking
open brackets and attempts to close the JSON at the latest point that still
parses, starting from the end and working backwards through the last twenty
candidate cut points. This exists because \texttt{inventory\_frame} and
\texttt{survey\_frame}, the two calls that ask for a list of objects, sometimes
hit \texttt{num\_predict} mid-list on a model that has started repeating itself,
and a reply with ten correct objects and one truncated one is more useful than
throwing the whole thing away. Every other call in this module keeps
\texttt{salvage=False}, since a truncated yes/no answer or a truncated bounding
box is not something a partial parse can safely recover.

Every stage built on \texttt{\_ask} is a separate, short, single-purpose prompt:
\texttt{check\_ambiguity}, \texttt{verify\_plan}, \texttt{describe\_frame},
\texttt{inventory\_frame}, \texttt{locate\_target}, \texttt{check\_condition} and
\texttt{survey\_frame} (the last two added for the conditional and surroundings
work, Sections~\ref{sec:conditional} and~\ref{ssec:survey}). None of them share a
prompt template; each states only what that one check needs to know. I did this
deliberately rather than writing one general-purpose ``look and answer'' function,
since a prompt that tries to serve several callers tends to grow extra fields
that most callers never fill in, which makes a wrong answer harder to trace back
to a cause.

\texttt{locate\_target} is the one with the most moving parts, since it is called
once per cycle of the free-roam approach (Section~\ref{ssec:freeroam}). It asks
for a bounding box, a confidence label, and, when \texttt{VLM\_DIST\_MAX\_CM}
permits it, a distance estimate in the units \texttt{VLM\_DIST\_MODE} selects:
either a number of metres or one of five coarse buckets. \texttt{distance\_cm\_from}
converts whichever the model returned into centimetres, or returns \texttt{None}
for any value outside \((0,\;\texttt{VLM\_DIST\_MAX\_CM}]\), including a
\texttt{NaN} that would otherwise fail every comparison silently. Every numeric
field from the model is treated this way: accepted only inside a sane range, and
\texttt{None} otherwise, because a model this size occasionally returns a
plausible-looking number for a question it cannot actually answer.

%% ===================================================================
\mysection{Clarification and mission orchestration}{Clarification and Mission Orchestration}
\label{sec:orchestration}

\mysubsection{Clarification manager}{Clarification Manager}
\label{sec:clarification_impl}

The clarification loop (Section~\ref{sec:clarification_design}) is implemented as
a small state machine, \texttt{DialogueSession}, and a module of handler
functions in \texttt{pipeline.py} that mutate it. \texttt{Phase} is a plain
enumeration (\texttt{CLARIFYING}, \texttt{PLANNING}, \texttt{VERIFYING},
\texttt{AWAITING\_CONFIRMATION}, \texttt{EXECUTING}, \texttt{REPORTING},
\texttt{CANCELLED}), and every \texttt{Turn} the rover and the person exchange is
appended to \texttt{session.turns}. \texttt{effective\_command} is the one
property everything downstream reads instead of the raw transcript: it returns
the resolved command once the vision model has rewritten one, the original
command once the session is capped, and otherwise the command with the
dialogue so far appended, in that order.

\texttt{\_advance\_to\_confirmation} is the function that drives one turn of the
loop. It calls \texttt{check\_ambiguity}, and if the reply says the command is
still ambiguous, it checks the proposed question against every question already
asked in this session, case-insensitively, before appending it. Without that
check, a model that cannot think of anything new to ask would repeat the same
question until the five-turn cap was reached, rather than giving up as soon as it
had nothing left to add. A session that reaches the cap, or runs out of new
questions, is marked \texttt{capped}, and \texttt{summarise\_plan}
(Section~\ref{sec:clarification_design}) opens the read-back with a warning in
that case.

Confirmation is a smaller state machine of its own, in \texttt{confirmation.py}.
\texttt{keyword\_pass} checks the reply against two phrase lists, English and
Afrikaans, with the reject list checked first; only if neither list matches does
\texttt{classify} fall back to a planner-model call, and an unparseable or
missing answer from that call is treated as \texttt{UNCLEAR} rather than guessed.
\texttt{DialogueSession} hands off to \texttt{MissionSession} once confirmation
succeeds, and everything from there on, including recovery and revision, is the
orchestration described in the rest of this section.

\mysubsection{The mission execution loop}{The Mission Execution Loop}
\label{ssec:mission_loop}

Once a plan is confirmed, \texttt{run\_mission} is the single function that
drives it, one step at a time, until the mission reaches a terminal phase or
runs out of steps. Each iteration reads \texttt{mission.cursor} to get the next
step, sets the status light to \texttt{executing}, and then checks whether the
step is conditional: if it carries a \texttt{when} tag of \texttt{yes} or
\texttt{no} and a branch answer has already been recorded from an earlier
\texttt{check} step, a step whose tag does not match the branch is recorded as
\texttt{skipped} with the branch it was skipped for, and the cursor moves on
without ever calling \texttt{execute\_step}. This is the piece of
Section~\ref{sec:conditional} that is actually code: the planner only decides
which steps carry a \texttt{when} tag, and the loop here is what honours it.

Which function actually runs the step depends on what it is and what the
rover can do. A \texttt{check} step goes to \texttt{\_run\_check}, which calls
the vision model's \texttt{check\_condition} or \texttt{survey\_frame} rather
than moving anything. An \texttt{approach} step goes to
\texttt{approach\_controller}, the free-roam hop loop
(Section~\ref{ssec:freeroam}), but only when \texttt{NAV\_MODE} is
\texttt{free} and the rover object actually exposes \texttt{pivot},
\texttt{hop}, \texttt{get\_frame} and \texttt{set\_gimbal}; everything else,
including every step on \texttt{VirtualRover} and every step on a
line-following drive, goes straight to \texttt{rover.execute\_step} on a
worker thread. A keyframe monitor is spawned alongside a \texttt{follow\_line}
step, but only when the rover object has a \texttt{get\_frame} method at all:
neither \texttt{SimulatedRover} nor \texttt{VirtualRover} implements one, so
the entire mid-drive check path is switched off for them automatically,
without an explicit mode flag anywhere, rather than behaving differently on
a model that happens not to have a camera to poll.

What happens after a step finishes depends on its result's \texttt{status}.
An \texttt{ok} result on a \texttt{check} step records the branch answer for
later conditional steps; on a \texttt{scan} step, or on an \texttt{observe} or
\texttt{report} step whose command asked about the surroundings, it triggers
the L5 survey (Section~\ref{ssec:survey}). A \texttt{blocked} result hands off
to the guidance loop (Section~\ref{ssec:guidance_loop}): if that loop stages a
new plan, the outer loop \texttt{continue}s without advancing the cursor,
and the next iteration sees \texttt{mission.pending\_swap} set and applies it
by replacing \texttt{active\_plan} and resetting the cursor to 0 before
reading the next step; if the guidance loop instead ends the mission, the
outer loop \texttt{break}s. A \texttt{halted} result breaks immediately, with
no guidance offered, since that status means the step itself decided the
mission cannot continue, which is a narrower condition than being blocked and
askable.

The loop's \texttt{finally} block is where every non-aborted mission writes
its report, whether it finished cleanly, ran out of guidance, or was halted
outright, because a run that ended badly is exactly the run whose report
matters most for the evaluation. The one path that skips this is an
\texttt{asyncio.CancelledError}, raised when the operator sends an explicit
stop: that path calls \texttt{rover.halt()} immediately, sets the phase to
\texttt{ABORTED}, and writes the deterministic half of the report straight to
disk with \texttt{spoken=False}, so stopping the rover is never delayed by a
planner call or a spoken summary it has no time to give.

\mysubsection{Failure recovery in code}{Failure Recovery in Code}
\label{ssec:guidance_loop}

A \texttt{blocked} result is handed to \texttt{\_guidance\_loop}, which is
where the recovery behaviour of Section~\ref{sec:recovery} is actually
implemented. If \texttt{GUIDANCE\_ENABLED} is off, it halts immediately with
a spoken ``I am blocked and have stopped.'' and nothing else happens; this
exists mainly for test batches that want a plain pass/fail on the first
failure rather than a multi-turn conversation. Otherwise it builds a question
from a fixed template keyed on the failure reason:
\texttt{target\_not\_found}, \texttt{obstacle}, \texttt{arrival\_contradicted},
and \texttt{cycle\_budget\_exhausted} each get their own wording, with
anything else falling back to a generic ``I could not complete the step.''
It then asks the question, up to \texttt{MAX\_GUIDANCE\_TURNS} times (3 by default).

Each turn waits for an answer through \texttt{\_wait\_for\_guidance}, which has
two modes. In normal operation it polls \texttt{mission.pending\_guidance}
every 200\,ms until \texttt{GUIDANCE\_TIMEOUT\_S} (120\,s by default) expires
or an answer arrives, which is how \texttt{handle\_guidance}, fed by either
transcribed audio or typed text from the dashboard, gets a reply into the
loop without the two functions sharing a direct call. In tests,
\texttt{mission.guidance\_provider} can instead be set to a plain synchronous
function, which lets a batch of L0–L5 runs script an answer for every
question the loop asks without any audio or WebSocket round-trip at all; a
provider that runs out of scripted answers returns \texttt{"stop"} rather than
raising, so an under-specified test script fails the mission the same way a
silent human would, not with a stack trace.

An answer that declines is checked with \texttt{\_declines}, not
\texttt{confirmation.keyword\_pass}, for the reason noted in
Section~\ref{sec:clarification_impl}: the confirmation reject list includes
words like ``no'' and ``wrong'' that show up constantly inside a perfectly
good guidance answer such as ``there is no box on the left, try the shelf''.
\texttt{\_declines} only triggers on an answer that opens with an explicit
stop word (\texttt{stop}, \texttt{cancel}, \texttt{abort}, \texttt{give up},
\texttt{never mind}, \texttt{forget it}, \texttt{quit}), or on a bare answer
of three words or fewer that \texttt{confirmation.keyword\_pass} alone calls a
reject. Anything longer that is not an explicit stop is assumed to be an
instruction and goes to \texttt{replan\_from\_state}
(Section~\ref{sec:planner}) along with the steps already done, the failed
step, the reason, and the guidance conversation so far. A plan that comes
back is staged as \texttt{pending\_swap} and the function returns \texttt{True};
a reply that produced no usable steps is logged and the loop asks again,
without spending one of the three turns on a parse failure that was not the
human's fault. Running out of turns, a timeout, or an explicit decline each
set a distinct \texttt{halt\_reason}
(\texttt{guidance\_timeout}/\texttt{guidance\_declined}/\texttt{guidance\_exhausted})
that ends up in the run log, so a batch of failed runs can be told apart by
why they failed rather than only that they did.

\mysubsection{Mid-mission revision}{Mid-mission Revision}
\label{ssec:revision_impl}

Every frame that arrives during a mission passes through \texttt{ingest\_frame},
which gates three tiers of increasingly expensive work so that the VLM is
not asked to look at every frame a camera happens to send. Tier 0 is free: a
32$\times$32 greyscale thumbnail and a Laplacian-variance blur score, computed
off the event loop's thread pool so a burst of frames cannot stall the step
loop. A frame below \texttt{PERCEPTION\_MIN\_SHARPNESS} (25.0) is dropped as
motion blur, and a frame that arrives before
\texttt{PERCEPTION\_MIN\_INTERVAL\_S} (6\,s) has passed, or whose thumbnail has
not moved more than \texttt{PERCEPTION\_CHANGE\_THRESHOLD} (0.06) from the last
one analysed, is dropped as ``nothing has changed''. Only a frame that passes
all three reaches tier 1, a single \texttt{describe\_frame} caption, which is
always recorded as an observation regardless of mission phase. Tier 2 only
fires while the mission is \texttt{EXECUTING}: it asks \texttt{revise\_plan}
whether the remaining steps still make sense against what was just seen, and
most captions change nothing, so most frames never reach this tier either.

A proposal from tier 2 is never applied on the model's say-so. It is handed
to \texttt{revision.classify}, a plain rule, not a model call, which compares
the confirmed plan's targets against the proposed plan's targets and the
action vocabulary each uses. If the proposed plan adds a target, drops a
confirmed target, or introduces an action outside a fixed whitelist of
navigation words (\texttt{move}, \texttt{approach}, \texttt{turn},
\texttt{scan}, and so on, with no world-touching action on that list by
construction), the change is \texttt{MATERIAL} and the mission moves to
\texttt{AWAITING\_REVISION\_CONFIRMATION}, speaks the delta through
\texttt{revision.summarise}, and waits for a yes before
\texttt{handle\_revision\_confirmation} applies it; an empty revised plan is
\texttt{BLOCKED} and ends the mission outright. Only when the goal and every
target are unchanged, and the only new actions are navigation ones, is the
change a \texttt{REROUTE}: that case is applied immediately and only
announced, never confirmed, since nothing about what the rover is trying to
do has actually changed, only the path it takes there. I kept this
classification out of the model deliberately, for the same reason
\texttt{confirmation.keyword\_pass} exists ahead of a model fallback: whether
a human must be consulted is not a decision to delegate to something that can
hallucinate its way to ``no change needed''.

%% ===================================================================
\mysection{Rover client and reactive layer}{Rover Client and Reactive Layer}
\label{sec:rover_client}

\mysubsection{Real hardware: PiRoverController}{Real Hardware}
\label{ssec:pi_rover}

\texttt{RoverController} is an abstract class with two required methods,
\texttt{execute\_step} and \texttt{halt}, and one optional one,
\texttt{status\_light}, which defaults to a no-op so the simulated rovers do not
need to implement lighting. Three classes satisfy it: \texttt{SimulatedRover}, a
timer that always reports success and exists only to exercise the mission loop
without any geometry; \texttt{VirtualRover}, described in
Section~\ref{ssec:virtual_rover}; and \texttt{PiRoverController}, which drives
the real TurboPi.

\texttt{PiRoverController} talks to the Pi over \texttt{roslibpy}, a Python
client for rosbridge. On construction it opens one \texttt{roslibpy.Ros}
connection to the rosbridge websocket, advertises a \texttt{/cmd\_vel} topic for
driving and a PWM-servo topic for the gimbal, and opens two ROS2 services,
\texttt{/line\_follow/cmd} and \texttt{/line\_follow/set\_running}, that control
the stock line-follower node. The sonar reading is a subscribed topic throttled
to a 100\,ms publish rate on the Pi side; each accepted reading increments a
sequence counter, \texttt{\_sonar\_seq}, which is how the obstacle check tells
``two consecutive readings below the threshold'' from ``the same stale value read
twice'', since the subscription is throttled but the control loop that reads it
is not.

Two threading events carry meanings that are easy to conflate and are kept
deliberately separate. \texttt{\_halted} means the controller has been told to
stop, through \texttt{halt()}, and is cleared only by the next
\texttt{execute\_step}; this is the flag Section~\ref{sec:recovery} explains the
free-roam loop must never trip by calling \texttt{halt()} on its own.
\texttt{\_target\_confirmed} means a mid-drive scan has confirmed the named
target is visible, and is reported by \texttt{\_follow\_line} as a successful
arrival rather than an abort, since it is a different event from being stopped.

The line-follower supervisor, \texttt{\_follow\_line}, owns exactly the three
things the stock node has no opinion about, as Section~\ref{ssec:line_follower}
describes: arrival through the node's own \texttt{crossroads\_stop} topic,
obstacle stopping through the polled sonar, and giving up through a wall-clock
timeout, since the node itself will drive forward indefinitely if the line is
lost. As of 11 October it also checks a new flag, \texttt{\_lf\_unavailable},
set whenever either ROS service call fails: if the node's services cannot be
reached at start-up, the supervisor waits 0.3\,s and then returns \texttt{blocked}
with the reason \texttt{line\_follow\_unavailable}, instead of reporting a drive
that never actually started as if it had run and simply found nothing.

\mysubsection{Virtual testing: VirtualRover}{Virtual Testing}
\label{ssec:virtual_rover}

\texttt{VirtualRover} is what every run without physical hardware actually
executes against, and it is built around one idea: \texttt{approach} is a
navigation skill, so it plans a route and follows it, while \texttt{move} and
\texttt{turn} are literal, so they are executed exactly as commanded with no
obstacle avoidance of their own. A failure during \texttt{approach} can
therefore only ever be a planning failure (no path, or a target the room does
not have), never a collision, while a collision during \texttt{move} or
\texttt{turn} means the chain of plan steps itself went somewhere it should
not have, which is a meaningfully different failure for the evaluation to
tell apart.

Both halves read from the same \texttt{OccupancyGrid}
(\texttt{nav.py}): every obstacle polygon in the room, plus the room boundary,
inflated by the rover's footprint clearance, a single value computed once from
the rover's length and width (18.7\,cm by 16.2\,cm by default) plus a safety
margin. \texttt{approach} plans an A* route over this grid with
\texttt{nav.plan\_path} and then walks it waypoint by waypoint; \texttt{move}
and \texttt{turn} instead sweep the straight segment they are about to travel
against the same grid with \texttt{nav.sweep\_segment} and stop at the last
free point if it is blocked. Using one grid for both was a deliberate choice:
a planner and a collision checker built from two different descriptions of the
room could disagree about what is free, and a route the planner considers
clear running straight into a wall the checker knows about would be a bug in
the test harness, not a finding worth reporting.

The room itself, a \texttt{RoomMap} loaded from JSON
(\texttt{room\_map.py}), is hand-typed, not estimated: coordinates for every
named landmark are measured off a sketch and a tape measure once per test
room, in metres, with a landmark's point being the standing spot a rover can
occupy to look at the thing, not the thing's own centroid. A target the
command names that the room map does not resolve is not treated as a crash;
it is logged as a hallucination event, which is exactly the RQ2 signal the
evaluation needs, and the step still returns a normal \texttt{blocked} result
with \texttt{reason: unresolved\_target}.

Every step also applies a configurable drift: Gaussian jitter on the achieved
position after a move, and on the achieved heading after a move or a turn,
scaled by \texttt{VIRTUAL\_DRIFT\_FRAC} and \texttt{VIRTUAL\_DRIFT\_DEG}. Both
are left at zero in the evaluation sweeps, so the numbers in
Chapter~\ref{chap:evaluation} are not inflated or deflated by injected noise;
the mechanism exists so the rover's open-loop drive can be exercised under
slip without needing the physical chassis for every run. Simulated and
wall-clock time are also kept apart on purpose: every action's duration is
first computed in simulated time from the linear speed and turn rate
configured for it, and only then divided by \texttt{VIRTUAL\_TIME\_SCALE} to
decide how long the process actually sleeps, so a sweep of many runs can be
played back faster than real time without changing the distances, speeds, or
timings the evaluation reports on.

Every run writes one JSON line per step to a trace file under
\texttt{Dashboard/brain/logs}, carrying the planned and achieved pose, the
step's status and reason, the running odometer, and, for an \texttt{approach}
step, the whole planned path, so a run can be plotted after the fact without
re-running it. \texttt{summary()} rolls a run's trace up into the numbers
RQ2 is built on: step count, how many were blocked, how many targets were
hallucinated and which ones, collisions, how often no path existed at all,
and the run's final pose.

%% ===================================================================
\mysection{Communication protocol and dashboard}{Communication Protocol and Dashboard}
\label{sec:protocol}

The brain is a FastAPI application. The two endpoints that matter for a mission,
\texttt{/ws/dialogue} and \texttt{/ws/execution}, are WebSockets carrying small
JSON frames rather than a request-response API, since both sides need to push
events the other did not ask for (a spoken prompt, a camera frame, a status
light). Figure~\ref{fig:protocol} shows the frame types each side sends.

\begin{figure}[htbp]
  \centering
  \begin{tikzpicture}[node distance=1.0cm]
    \node[hw] (pi) at (0,0)   {Pi client};
    \node[blk] (br) at (8,0)  {Brain (FastAPI)};
    \draw[arr] ([yshift=0.9cm]pi.east) -- node[lbl, above]{start, audio\_chunk, confirmation\_audio}([yshift=0.9cm]br.west);
    \draw[arr] ([yshift=0.3cm]br.west) -- node[lbl, above]{session, speak, plan\_ready, execute, cancelled}([yshift=0.3cm]pi.east);
    \draw[arr] ([yshift=-0.3cm]pi.east) -- node[lbl, below]{begin, frame, revision\_audio}([yshift=-0.3cm]br.west);
    \draw[arr] ([yshift=-0.9cm]br.west) -- node[lbl, below]{step\_started, check, awaiting\_guidance, report}([yshift=-0.9cm]pi.east);
  \end{tikzpicture}
  \caption{\texttt{/ws/dialogue} (top pair) carries everything up to confirmation; \texttt{/ws/execution} (bottom pair) carries the mission. One frame per complete utterance, not a stream.}
  \label{fig:protocol}
\end{figure}

One frame carries one complete utterance or one complete camera frame, never a
stream of partial audio, for the same reason Section~\ref{sec:stt} gives:
Whisper pads every clip regardless, so sending a clip in pieces would add
reassembly complexity for no latency gain. Both connections are opened with
\texttt{max\_size=None}. This was not the original setting; a long spoken plan
read-back, a \texttt{speak} event carrying the base64 WAV Section~\ref{ssec:tts}
describes, exceeded the websockets library's 1\,MiB default and closed the
connection with code 1009, which I only found by reproducing a failure that
happened on longer plans and not shorter ones.

\texttt{/scene} accepts an uploaded walkthrough video and returns a scene id;
\texttt{/command} and \texttt{/transcribe} exist for one-shot testing outside a
WebSocket session; \texttt{/logs} and \texttt{/logs/\{session\_id\}} serve the run
log listing and detail that the dashboard's Logs page reads
(Section~\ref{sec:logging}); and \texttt{/logs/\{session\_id\}/flag} is a small
\texttt{PUT} endpoint, added alongside the L0 to L5 test batches, that lets an
operator mark a run good, bad or for review with a note, stored in a separate
\texttt{run\_flags.json} rather than inside the run log itself, so flagging a run
never touches the record the run itself wrote.

The dashboard is a React and TypeScript single-page app (Vite, the components
under \texttt{ui/src/components}) that is a thin client over these same
endpoints: a \texttt{DialoguePanel} and \texttt{ExecutionPanel} for driving and
watching a live session, a \texttt{GimbalControl} and \texttt{RoverMap} for
manual testing, and a \texttt{LogsPage} for reading back past runs and setting
the flag above. It holds no logic of its own that the evaluation depends on; it
exists so a run can be started and watched without a terminal, and every number
Chapter~\ref{chap:evaluation} reports comes from the JSON logs, not from the UI.

The dashboard itself is not part of the project's contribution, so I leaned on
an AI assistant to build most of its React components rather than writing them
by hand. It is a debugging and demonstration tool, nothing more: something to
let me upload a room walkthrough video, watch a dialogue and a mission happen
panel by panel, and drive the rover manually, without needing a terminal and a
stack of \texttt{curl} commands to check that the pipeline underneath was
actually working end to end. Using AI for this piece made sense because the
time saved on React boilerplate went into the planner, the vision model and
the mission logic instead, which is where the actual research question lives.

%% ===================================================================
\mysection{Logging and measurement instrumentation}{Logging and Measurement Instrumentation}
\label{sec:logging}

Two logs exist side by side and serve different purposes. \texttt{brain.log} is
a conventional application log, written through Python's \texttt{logging} module
via \texttt{log\_setup.get\_logger}: every subsystem gets a named logger
(\texttt{brain.planner}, \texttt{brain.vlm}, \texttt{brain.mission}, and so on)
under one root handler, so a single file interleaves every call in time order
with full prompts and raw model output at debug level, while the console only
shows info level and above. This log is for debugging a run, not for the
evaluation; it is not structured and I do not parse it programmatically.

The run log, one JSON file per mission under \texttt{logs/mission\_\textit{id}.json},
is the one the evaluation is built on, and \texttt{report.py} is what assembles
it. It is written once, at the end of a mission, from everything the mission
session accumulated: the plan, every check and its result, the guidance log,
the light-state timeline, the model tag for both roles and whether the run
was local or cloud, and a tagged \texttt{test\_level} when one is set, so a
batch of L0 to L5 runs can be grouped without reading every command by hand.
\texttt{report.py}'s \texttt{\_outcome\_class} is what turns the raw final
phase into one of \texttt{completed}, \texttt{completed\_after\_guidance},
\texttt{halted\_after\_guidance}, \texttt{halted} or \texttt{aborted}: any
mission that needed at least one guidance turn is never reported as a plain
\texttt{completed}, even if it finished, because Section~\ref{sec:recovery}'s
distinction between a clean run and a recovered one is exactly what RQ1 needs
to see. \texttt{\_reached\_after\_guidance} narrows this further for the
\texttt{recovered} flag: it only counts as recovered if an \texttt{approach}
step actually succeeded somewhere after the last \texttt{blocked} step, not
merely if the mission eventually ended without halting.

Every model call, in both the planner and the vision model, also writes one
entry to a smaller usage ledger (\texttt{providers/usage.py}), holding latency,
token counts and cost when the provider reports them; this ledger is reset at
the start of every \texttt{/ws/dialogue} session and is what \texttt{report.py}'s
\texttt{usage\_summary} sums over, since only one mission runs per process at a
time.

A separate, offline tool, \texttt{tools/context\_budget.py}, measures something
the run log cannot: how many tokens a given prompt would cost before it is ever
sent. It imports \texttt{planner.py} and \texttt{vlm.py} directly, forces
\texttt{ROVER=virtual}, and asks Ollama for the \texttt{prompt\_eval\_count} on a
constructed prompt with \texttt{num\_predict=1}, which counts the tokens without
paying for a real generation. It measures the fixed instructions, the room
digest, one clarification turn and one recovery turn separately, so the
measurements in Chapter~\ref{chap:evaluation} can say how many turns of each kind
fit inside the 8192-token window before it becomes the limiting factor
(Section~\ref{ssec:split}), rather than only reporting that a run eventually
failed.
